\documentclass[10pt,a4paper]{amsart}
\usepackage[margin=19mm]{geometry}
\usepackage{amsmath,amssymb,mathtools}
\usepackage[scr=rsfs,cal=euler]{mathalfa}
\usepackage{xfrac}
\usepackage[unicode,hidelinks,bookmarks=false]{hyperref}
\newcommand{\ie}{i.e.\ }
\newcommand{\cf}{cf.\ }
\newcommand{\resp}{respectively\ }
\newcommand{\Subsection}{\subsection}
\providecommand{\nobreakdash}{\nobreak\hbox{-}}
\setlength{\emergencystretch}{2em}
\multlinegap0pt
\usepackage{amssymb}
\usepackage{float}
\usepackage{xcolor}
\newtheorem{lemma}{Lemma}
\title{On Hales-Jewett and Related Numbers}
\author{Nathan Conlon, Nayda Farnsworth,, Aaron Robertson}
\date{Source: arXiv:2607.18111v1 (CC BY 4.0)}
\begin{document}
\maketitle

\noindent\emph{Source and licence.} On Hales-Jewett and Related Numbers. Version arXiv:2607.18111v1 (CC BY 4.0), distributed by arXiv under the Creative Commons Attribution 4.0 International licence, http://creativecommons.org/licenses/by/4.0/, as recorded in the arXiv licence metadata for this exact version and shown on the abstract page https://arxiv.org/abs/2607.18111v1. Full licence notice: \texttt{licenses/arxiv-2607.18111-CC-BY-4.0.txt}.\par\smallskip

\noindent\textit{Editorial scope.} Counted unit: the article's lemma lem1 (source0.tex lines 1075--1076). The article's complete original proof of it is delivered with it (source0.tex lines 1078--1112, one \texttt{proof} environment of 35 lines). No mathematics is written, repaired or re-derived: the statement and the proof are the article's own lines, in the article's own notation, and no retained line is substituted or re-flowed.  No further source-local context is needed: every cross-reference of every retained line resolves inside the two delivered spans.  Nothing was omitted from any retained span, so the package carries no omission record.  No retained line contains a citation, so the wrapper carries no bibliography and no bibliography entry.  No retained line contains a figure environment, an image-inclusion command or drawing code, so no image is delivered and the package declares no asset dependency.  Editorial formatting only, in addition: an AMS \texttt{amsart} preamble that restates the article's own declarations and macros used by the retained lines, namely the environments \texttt{lemma} on separate counters, together with the standard packages those lines need.  The retained environments print in this document as Lemma 1 in order of appearance, whereas the article prints the same items with the numbering of its own full document.  The complete native source of the article as distributed in the arXiv e-print for this exact version is bundled unchanged as \texttt{sources/205626/source0.tex}.
\begin{lemma}\label{lem1} For all $k\in\mathbb{Z}^+$ we have $a(k) = p(k-1)+2$, where $p$ is the least prime factor of $k$.
\end{lemma}

\begin{proof} We start by providing a valid 2-coloring of $[1,p(k-1)+1]$.
Consider the 2-coloring of $[1,p(k-1)+1]$ defined by coloring the congruence class $1$ modulo $k$
blue and all other integers red.  Since this coloring has no two   blue
integers with difference less than $k$ we need only show that this does not admit a red $k$-ap.  

Let
$a,a+d,\dots,a+(k-1)d$ be an arbitrary $k$-ap.  So that $a+(k-1)d \leq p(k-1)+1$ we must have $d < p$ or $d=p$ and $a=1$.
In the latter case, we see that $a$ is blue so that the $k$-ap is not red.  Hence, we may now assume $d<p$
so that $\gcd(d,k)=1$.  Hence, $xd \equiv 1-a \pmod{k}$ has a solution, say $x=j$,
with $0 \leq j \leq k-1$.  This means that $a+jd \equiv 1 \pmod{k}$ so that it has color blue.  Hence, for $d<p$ there
is no red $k$-ap.  

As all cases for the value of $d$ have been exhausted, we have a valid 2-coloring of $[1,p(k-1)+1]$.
This completes the lower bound.

For the upper bound, assume, for a contradication, that $\chi$ is a valid coloring of $[1,p(k-1)+2]$; that is, 
there is no red $k$-ap and no two blue integers
with difference less than $k$.  

Since any $k$ consecutive integers form a $k$-ap, every set of
$k$ consecutive integers must contain a blue integer.  
Hence, since every two consecutive blue
integers have a difference of at least $k$, we conclude that for some $m \in \{1,2,\dots,k\}$ the congruence
class $m$ modulo $k$ consists of only blue integers and all other integers are red.  

We now turn our attention to $m$.  Reduce $m$ modulo $p$.  Note that $m$ cannot be both
congruent to $1$ and $2$ modulo $p$.  Hence, one of $\{1+ip: 0 \leq i \leq k-1\}$
and $\{2+ip: 0 \leq i \leq k-1\}$ does not intersect the congruence class $m$ modulo $p$, and,
consequently, does not intersect the congruence class $m$ modulo $k$.
Since the largest element in these two $k$-term arithmetic progressions is $p(k-1)+2$ and
one of them must be entirely red, we have a contradiction to $\chi$ being valid.
This completes the upper bound.

Since the upper and lower bounds agree, the lemma holds.
\end{proof}

\end{document}
